\section{De la saturación de los puntos de referencia a los Agentes multimodales: por qué se reescribió el espacio del problema}

\subsection{Hilo conductor de la charla y juicio central}
Esta charla la dio Caiming Xiong, de Salesforce AI Research, y su hilo conductor es muy claro: los puntos de referencia tradicionales de tareas de texto se acercaron a la saturación en los últimos años, y el desafío de la siguiente etapa ya no es «responder preguntas fuera de línea», sino «completar tareas de manera continua en un entorno digital real». Esta definición desplaza el objeto de estudio de tareas de PLN estáticas hacia sistemas de Agentes interactivos, de varios pasos y que operan a través de aplicaciones.

\begin{importantbox}{De «saber responder» a «saber hacer»}
En el paradigma de los Agentes multimodales, el valor de un modelo ya no se mide solo por si una frase que produce es correcta, sino por si puede completar de manera estable un flujo de trabajo real: leer la interfaz, entender el estado, planear la acción, ejecutar la operación, verificar el resultado y volver a iterar. Este ciclo cerrado es el verdadero umbral en un entorno de producción.
\end{importantbox}

\subsection{Por qué «multimodal» ya no es opcional}
El expositor define al Agente multimodal como un sistema Vision-Language-Action (VLA). La razón es directa: el flujo de trabajo real no es una API de puro texto, sino sobre todo interfaces gráficas, páginas web, software de escritorio, interfaces móviles e incluso interacción con periféricos. Solo si el Agente cuenta a la vez con comprensión visual, razonamiento en lenguaje y generación de acciones puede llegar a hacerse cargo de un proceso complejo.

\begin{knowledgebox}{Comparación de entradas y salidas del Agente multimodal}
\begin{itemize}
    \item La entrada pasa de ser «texto de un solo turno» a ser «instrucción del usuario + captura de pantalla + árbol de accesibilidad + estado del entorno»;
    \item la salida pasa de ser «respuesta en lenguaje natural» a ser «secuencia de acciones ejecutables»;
    \item la evaluación pasa de ser «exactitud» a ser «tasa de finalización de la tarea, calidad de la trayectoria, eficiencia de los pasos y robustez».
\end{itemize}
\end{knowledgebox}

\begin{figure}[H]
\centering
\includegraphics[width=0.88\textwidth]{frame-01-benchmark.jpg}
\caption{Video, aproximadamente 00:01:30: el expositor muestra que los puntos de referencia recientes se saturan con rapidez, y a partir de ahí plantea que la evaluación de Agentes debe pasar de un banco de preguntas fuera de línea a un entorno real}
\end{figure}

\begin{warningbox}{Malentendido frecuente: confundir «saber usar herramientas» con «Agente listo para producción»}
Un éxito en una demostración no significa que el sistema esté listo para salir a producción. Un Agente de nivel de producción se preocupa por la estabilidad a largo plazo: recuperación de errores, manejo de interrupciones de tareas, consistencia del estado entre aplicaciones, límites de permisos y capacidad de auditoría. Ninguno de estos aspectos puede cubrirse con la puntuación de un solo benchmark.
\end{warningbox}

\subsection{Resumen de la sección}
El foco de la investigación sobre Agentes multimodales pasó de la «capacidad de responder» del modelo a la «capacidad de interactuar con el entorno». Este giro determina que, de ahora en adelante, deban construirse a la vez tres capas de base: un benchmark de entorno real, un pipeline de datos escalable y una arquitectura de modelo ejecutable.
